\section{2026-05-12}

\subsection{Real Projective Space}

\begin{definition}
The real projective space $\mathbb{RP}^n$ is defined by
\[
  \mathbb{RP}^n
  :=
  \left(\RR^{n+1} \setminus \{0\}\right)/\sim,
\]
where
\[
  x \sim y
  \iff
  x = \lambda y
  \text{ for some } \lambda \in \RR \setminus \{0\}.
\]
Equivalently,
\[
  \mathbb{RP}^n
  =
  \left\{
    L \subset \RR^{n+1}
    \mid
    L \text{ is a line through the origin}
  \right\}.
\]
\end{definition}

Let
\[
  \pi : \RR^{n+1} \setminus \{0\} \to \mathbb{RP}^n
\]
be the quotient map. For $x \in \RR^{n+1} \setminus \{0\}$, we write
\[
  \pi(x) = [x].
\]
The quotient topology on $\mathbb{RP}^n$ is defined as follows:
\[
  A \subset \mathbb{RP}^n \text{ is open}
  \iff
  \pi^{-1}(A) \text{ is open in } \RR^{n+1} \setminus \{0\}.
\]

Equivalently, one can also describe $\mathbb{RP}^n$ as
\[
  S^n/(x \sim -x).
\]
The quotient map is
\[
  \pi : S^n \to \mathbb{RP}^n,
  \qquad
  x \mapsto [x].
\]
This gives the same quotient topology on $\mathbb{RP}^n$.

\begin{definition}
A point of $\mathbb{RP}^n$ is written in homogeneous coordinates as
\[
  [x^1 : x^2 : \cdots : x^{n+1}].
\]
By definition,
\[
  [x^1 : x^2 : \cdots : x^{n+1}]
  =
  [\lambda x^1 : \lambda x^2 : \cdots : \lambda x^{n+1}]
\]
for every $\lambda \in \RR \setminus \{0\}$.
\end{definition}

For each $i = 1,\dots,n+1$, define
\[
  U_i
  :=
  \left\{
    [x^1 : \cdots : x^{n+1}]
    \in \mathbb{RP}^n
    \mid
    x^i \neq 0
  \right\}.
\]

The set $U_i$ is open in $\mathbb{RP}^n$. Indeed,
\[
  \pi^{-1}(U_i)
  =
  \left\{
    x \in \RR^{n+1} \setminus \{0\}
    \mid
    x^i \neq 0
  \right\},
\]
which is open in $\RR^{n+1} \setminus \{0\}$. Hence $U_i$ is open by the definition of the quotient topology.

Moreover, the sets $U_1,\dots,U_{n+1}$ cover $\mathbb{RP}^n$. Indeed, for every
\[
  [x^1 : \cdots : x^{n+1}] \in \mathbb{RP}^n,
\]
not all coordinates $x^1,\dots,x^{n+1}$ are zero. Hence at least one coordinate $x^i$ is nonzero, so the point lies in $U_i$.

\begin{definition}
For each $i = 1,\dots,n+1$, define
\[
  \varphi_i : U_i \to \RR^n
\]
by
\[
  \varphi_i\left([x^1 : \cdots : x^{n+1}]\right)
  =
  \left(
    \frac{x^1}{x^i},
    \dots,
    \frac{x^{i-1}}{x^i},
    \frac{x^{i+1}}{x^i},
    \dots,
    \frac{x^{n+1}}{x^i}
  \right).
\]
\end{definition}

The map $\varphi_i$ is well-defined. Indeed, if
\[
  [x^1 : \cdots : x^{n+1}]
  =
  [\lambda x^1 : \cdots : \lambda x^{n+1}],
  \qquad \lambda \neq 0,
\]
then for every $j \neq i$,
\[
  \frac{\lambda x^j}{\lambda x^i}
  =
  \frac{x^j}{x^i}.
\]

The inverse map is
\[
  \varphi_i^{-1}(u^1,\dots,u^n)
  =
  [u^1 : \cdots : u^{i-1} : 1 : u^i : \cdots : u^n].
\]
Therefore $\varphi_i$ is a bijection. Moreover, both $\varphi_i$ and
$\varphi_i^{-1}$ are continuous, so $\varphi_i$ is a homeomorphism.

Hence
\[
  (U_i,\varphi_i)
\]
is a coordinate chart on $\mathbb{RP}^n$.

\begin{theorem}
The collection
\[
  \mathcal A
  =
  \{(U_i,\varphi_i) \mid i = 1,\dots,n+1\}
\]
is a smooth atlas on $\mathbb{RP}^n$.
\end{theorem}

\begin{proof}
The sets $U_1,\dots,U_{n+1}$ cover $\mathbb{RP}^n$, so it remains to check that the charts are smoothly compatible.

Let $i \neq j$. On the overlap $U_i \cap U_j$, both $x^i$ and $x^j$ are nonzero. Consider the transition map
\[
  \varphi_j \circ \varphi_i^{-1}.
\]
In the $i$-th chart, a point of $U_i$ has the form
\[
  [x^1 : \cdots : x^{n+1}]
  =
  \left[
    \frac{x^1}{x^i}
    :
    \cdots
    :
    \frac{x^{i-1}}{x^i}
    :
    1
    :
    \frac{x^{i+1}}{x^i}
    :
    \cdots
    :
    \frac{x^{n+1}}{x^i}
  \right].
\]
Passing to the $j$-th chart means dividing every coordinate by the $j$-th coordinate. Hence the transition map is given by rational functions whose denominators are nonzero on the overlap.

Therefore
\[
  \varphi_j \circ \varphi_i^{-1}
\]
is smooth. Similarly,
\[
  \varphi_i \circ \varphi_j^{-1}
\]
is smooth. Thus the charts are smoothly compatible.

Therefore $\mathcal A$ is a smooth atlas on $\mathbb{RP}^n$.
\end{proof}


\subsection{Polar Coordinates and Transition Maps}

\begin{example}
Let $M=\RR^2$.

Consider polar coordinates on
\[
  U=\RR^2\setminus \left(\RR_{\ge 0}\times \{0\}\right).
\]
For $z\in U$, write $z=(z^1,z^2)$. Define
\[
  \varphi:U\to \RR^2
\]
by
\[
  \varphi(z)
  =
  \bigl(r(z),\theta(z)\bigr)
  =
  \bigl(\varphi^1(z),\varphi^2(z)\bigr)
  =
  \bigl(|z|,\arg z\bigr),
\]
where
\[
  r,\theta:U\to \RR
\]
are the coordinate functions.

Here the image of $\varphi$ is
\[
  \varphi(U)=\RR_+\times (0,2\pi).
\]
Thus the polar coordinate chart is
\[
  (U,\varphi).
\]
\end{example}

\begin{example}
Now consider the Euclidean coordinates on
\[
  V=\RR^2.
\]
Define
\[
  \psi:V\to \RR^2
\]
by
\[
  \psi(z)
  =
  \bigl(x(z),y(z)\bigr)
  =
  \bigl(\psi^1(z),\psi^2(z)\bigr)
  =
  (z^1,z^2),
\]
where
\[
  x,y:\RR^2\to \RR
\]
are the coordinate functions.

Thus the Euclidean coordinate chart is
\[
  (V,\psi).
\]
\end{example}

The transition map from polar coordinates to Euclidean coordinates is
\[
  \psi\circ \varphi^{-1}:\RR_+\times (0,2\pi)
  \to
  \RR^2\setminus \left(\RR_{\ge 0}\times \{0\}\right).
\]
Explicitly,
\[
  (\psi\circ \varphi^{-1})(\widehat r,\widehat\theta)
  =
  \bigl(\widehat r\cos \widehat\theta,
  \widehat r\sin \widehat\theta\bigr).
\]

Since the component functions
\[
  (\widehat r,\widehat\theta)
  \mapsto
  \widehat r\cos \widehat\theta
  \quad\text{and}\quad
  (\widehat r,\widehat\theta)
  \mapsto
  \widehat r\sin \widehat\theta
\]
are smooth, the transition map
\[
  \psi\circ \varphi^{-1}
\]
is smooth.

\subsection{Smooth Compatibility}

\begin{definition}
Let $(U,\varphi)$ and $(V,\psi)$ be two charts on a topological manifold $M$.
We say that $(U,\varphi)$ and $(V,\psi)$ are \vocab{smoothly compatible} if both transition maps
\[
  \psi\circ \varphi^{-1} :
  \varphi(U\cap V) \to \psi(U\cap V)
\]
and
\[
  \varphi\circ \psi^{-1} :
  \psi(U\cap V) \to \varphi(U\cap V)
\]
are smooth maps between open subsets of Euclidean spaces.
\end{definition}

\begin{remark}
At this stage, we have not yet defined what it means for a map between manifolds to be smooth.
Therefore we do not say that a chart itself is smooth. Instead, we say that two charts are smoothly compatible.
\end{remark}

\begin{remark}
If two charts $(U,\varphi)$ and $(V,\psi)$ are smoothly compatible, then the transition maps
\[
  \psi\circ \varphi^{-1}
  \quad\text{and}\quad
  \varphi\circ \psi^{-1}
\]
are diffeomorphisms between open subsets of Euclidean spaces.
\end{remark}

\subsection{Smooth Atlases and Maximal Atlases}

\begin{definition}
An \vocab{atlas} $\mathcal A$ of an $n$-dimensional topological manifold $M$ is a collection of charts
\[
  \mathcal A=\{(U_\alpha,\varphi_\alpha)\}_{\alpha\in I}
\]
such that
\[
  M=\bigcup_{\alpha\in I}U_\alpha.
\]
\end{definition}

\begin{definition}
An atlas $\mathcal A$ is called a \vocab{smooth atlas} if all charts in $\mathcal A$ are pairwise smoothly compatible.
\end{definition}

\begin{definition}
A smooth atlas $\mathcal A$ is called a \vocab{maximal smooth atlas} if there is no smooth atlas $\mathcal A'$ such that
\[
  \mathcal A \subsetneq \mathcal A'.
\]
Equivalently, $\mathcal A$ is maximal if every chart smoothly compatible with all charts in $\mathcal A$ is already contained in $\mathcal A$.
\end{definition}

\begin{theorem}
Every smooth atlas $\mathcal A$ of $M$ is contained in a unique maximal smooth atlas.
\end{theorem}

\begin{proof}
We first prove existence.

Define
\[
  \overline{\mathcal A}
  :=
  \left\{
    (U,\varphi)
    \ \middle|\
    \begin{array}{l}
      (U,\varphi) \text{ is a chart of } M,\\
      (U,\varphi) \text{ is smoothly compatible with every chart in } \mathcal A
    \end{array}
  \right\}.
\]
Clearly,
\[
  \mathcal A\subset \overline{\mathcal A}.
\]

\begin{claim*}
$\overline{\mathcal A}$ is a smooth atlas.
\end{claim*}

\begin{proof}
Since $\mathcal A\subset \overline{\mathcal A}$ and $\mathcal A$ covers $M$, the collection $\overline{\mathcal A}$ also covers $M$.

It remains to prove that any two charts in $\overline{\mathcal A}$ are smoothly compatible.

Let
\[
  (U_1,\varphi_1),(U_2,\varphi_2)\in \overline{\mathcal A}.
\]
If
\[
  U_1\cap U_2=\varnothing,
\]
then the transition maps have empty domain, so the two charts are smoothly compatible vacuously.

Now suppose
\[
  U_1\cap U_2\neq \varnothing.
\]
For each point
\[
  p\in U_1\cap U_2,
\]
since $\mathcal A$ is an atlas, there exists a chart
\[
  (V,\psi)\in \mathcal A
\]
such that
\[
  p\in V.
\]
By definition of $\overline{\mathcal A}$, both $(U_1,\varphi_1)$ and $(U_2,\varphi_2)$ are smoothly compatible with $(V,\psi)$.

Therefore the transition maps
\[
  \psi\circ \varphi_1^{-1}
  \quad\text{and}\quad
  \varphi_2\circ \psi^{-1}
\]
are smooth on their domains. On the open set
\[
  \varphi_1(U_1\cap U_2\cap V),
\]
we have
\[
  \varphi_2\circ \varphi_1^{-1}
  =
  \left(\varphi_2\circ \psi^{-1}\right)
  \circ
  \left(\psi\circ \varphi_1^{-1}\right).
\]
Hence $\varphi_2\circ \varphi_1^{-1}$ is smooth near $\varphi_1(p)$.

Since this holds for every $p\in U_1\cap U_2$, the transition map
\[
  \varphi_2\circ \varphi_1^{-1}:
  \varphi_1(U_1\cap U_2)\to \varphi_2(U_1\cap U_2)
\]
is smooth.

Similarly,
\[
  \varphi_1\circ \varphi_2^{-1}
\]
is smooth. Therefore $(U_1,\varphi_1)$ and $(U_2,\varphi_2)$ are smoothly compatible.

Thus $\overline{\mathcal A}$ is a smooth atlas.
\end{proof}

% preamble:
% \usepackage{tikz}
% \usetikzlibrary{arrows.meta}

\begin{figure}[htbp]
\centering
\begin{tikzpicture}[
  >=Stealth,
  every node/.style={font=\small},
  map/.style={->, thick},
  smoothmap/.style={->, thick, red!70!black},
  outline/.style={thick},
  coordbox/.style={
    draw,
    thick,
    rounded corners=2pt,
    minimum width=4.2cm,
    minimum height=2.8cm
  }
]

% =====================================================
% Top: manifold M
% =====================================================

\node at (0,8.8) {$M$};

% M
\draw[outline]
  (0,7.5) ellipse (5.5 and 1.35);

% -----------------------------------------------------
% Exact shading of U_1 \cap U_2
% -----------------------------------------------------
\begin{scope}
  \clip (-1.4,7.5) ellipse (2.8 and 0.95); % U_1
  \clip ( 1.4,7.5) ellipse (2.8 and 0.95); % U_2
  \fill[orange!35, opacity=0.9] (-6.2,6.0) rectangle (6.2,9.0);
\end{scope}

% -----------------------------------------------------
% Exact shading of U_1 \cap U_2 \cap V
% -----------------------------------------------------
\begin{scope}
  \clip (-1.4,7.5) ellipse (2.8 and 0.95); % U_1
  \clip ( 1.4,7.5) ellipse (2.8 and 0.95); % U_2
  \clip ( 0.0,7.25) ellipse (1.35 and 0.60); % V
  \fill[red!55, opacity=0.75] (-6.2,6.0) rectangle (6.2,9.0);
\end{scope}

% outlines
\draw[outline, blue!70!black]
  (-1.4,7.5) ellipse (2.8 and 0.95);
\node[blue!70!black] at (-3.25,8.2) {$U_1$};

\draw[outline, red!70!black]
  (1.4,7.5) ellipse (2.8 and 0.95);
\node[red!70!black] at (3.25,8.2) {$U_2$};

\draw[outline, green!50!black]
  (0.0,7.25) ellipse (1.35 and 0.60);
\node[green!50!black] at (-1.35,6.45) {$V$};

\node at (0,8.0) {$U_1\cap U_2$};
\node[red!70!black] at (2.5,6.85) {$U_1\cap U_2\cap V$};

% point p
\fill (0,7.25) circle (1.8pt);
\node[above] at (0,7.30) {$p$};

% =====================================================
% Bottom: coordinate boxes
% =====================================================

% Left box
\node[coordbox] (L) at (-4.8,3.3) {};
\node at (-4.8,4.9) {$\RR^n$};
\node at (-4.8,4.4) {$\varphi_1(U_1)$};

\draw[outline] (-4.8,3.35) ellipse (1.40 and 0.56);
\begin{scope}
  \clip (-4.8,3.35) ellipse (1.40 and 0.56);
  \fill[orange!35, opacity=0.9] (-4.55,3.40) ellipse (0.82 and 0.30);
  \fill[red!55, opacity=0.75] (-4.95,3.18) ellipse (0.58 and 0.18);
\end{scope}
\node at (-4.45,3.90) {$\varphi_1(U_1\cap U_2)$};
\node[red!70!black] at (-4.8,2.25) {$\varphi_1(U_1\cap U_2\cap V)$};

\fill (-4.95,3.18) circle (1.4pt);
\node[below] at (-4.95,3.08) {$\varphi_1(p)$};

% Middle box
\node[coordbox] (C) at (0,3.3) {};
\node at (0,4.9) {$\RR^n$};
\node at (0,4.4) {$\psi(V)$};

\draw[outline] (0,3.35) ellipse (1.30 and 0.52);
\begin{scope}
  \clip (0,3.35) ellipse (1.30 and 0.52);
  \fill[red!55, opacity=0.75] (0,3.18) ellipse (0.62 and 0.19);
\end{scope}
\node[red!70!black] at (0,2.25) {$\psi(U_1\cap U_2\cap V)$};

\fill (0,3.18) circle (1.4pt);
\node[below] at (0,3.08) {$\psi(p)$};

% Right box
\node[coordbox] (R) at (4.8,3.3) {};
\node at (4.8,4.9) {$\RR^n$};
\node at (4.8,4.4) {$\varphi_2(U_2)$};

\draw[outline] (4.8,3.35) ellipse (1.40 and 0.56);
\begin{scope}
  \clip (4.8,3.35) ellipse (1.40 and 0.56);
  \fill[orange!35, opacity=0.9] (4.55,3.40) ellipse (0.82 and 0.30);
  \fill[red!55, opacity=0.75] (4.45,3.18) ellipse (0.58 and 0.18);
\end{scope}
\node at (4.50,3.90) {$\varphi_2(U_1\cap U_2)$};
\node[red!70!black] at (4.8,2.25) {$\varphi_2(U_1\cap U_2\cap V)$};

\fill (4.45,3.18) circle (1.4pt);
\node[below] at (4.45,3.08) {$\varphi_2(p)$};

% =====================================================
% Chart arrows from M
% =====================================================

\draw[map]
  (-2.9,6.6) to[out=-120,in=90] (-5.4,4.65);
\node at (-5.55,5.85) {$\varphi_1$};

\draw[map, green!50!black]
  (0,6.6) -- (0,4.65);
\node[green!50!black] at (0.35,5.55) {$\psi$};

\draw[map]
  (2.9,6.6) to[out=-60,in=90] (5.4,4.65);
\node at (5.55,5.85) {$\varphi_2$};

% =====================================================
% Transition maps
% =====================================================

\draw[smoothmap] (-2.55,3.35) -- (-2.05,3.35);
\node[red!70!black, above] at (-2.30,3.58) {$\psi\circ\varphi_1^{-1}$};
\node[red!70!black, below] at (-2.30,3.10) {smooth};

\draw[smoothmap] (2.05,3.35) -- (2.55,3.35);
\node[red!70!black, above] at (2.30,3.58) {$\varphi_2\circ\psi^{-1}$};
\node[red!70!black, below] at (2.30,3.10) {smooth};

% =====================================================
% Bottom formula
% =====================================================

\node[align=center] at (0,0.85)
{On $\varphi_1(U_1\cap U_2\cap V)$,\\[0.4em]
$\displaystyle
  \varphi_2\circ\varphi_1^{-1}
  =
  (\varphi_2\circ\psi^{-1})\circ(\psi\circ\varphi_1^{-1})$\\[0.4em]
Hence $\varphi_2\circ\varphi_1^{-1}$ is smooth near $\varphi_1(p)$.
};

\end{tikzpicture}
\caption{The local factorization used in the proof of the first claim of Theorem 1.13.}
\end{figure}

Now we prove maximality.

\begin{claim*}
If $\mathcal A'$ is a smooth atlas with $\mathcal A\subset \mathcal A'$, then $\mathcal A'\subset \overline{\mathcal A}$.
\end{claim*}

\begin{proof}
Let
\[
  (U',\varphi')\in \mathcal A'.
\]
Since $\mathcal A'$ is a smooth atlas, $(U',\varphi')$ is smoothly compatible with every chart in $\mathcal A'$.
Since
\[
  \mathcal A\subset \mathcal A',
\]
it follows that $(U',\varphi')$ is smoothly compatible with every chart in $\mathcal A$.

Therefore, by definition of $\overline{\mathcal A}$,
\[
  (U',\varphi')\in \overline{\mathcal A}.
\]
Hence
\[
  \mathcal A'\subset \overline{\mathcal A}.
\]
\end{proof}

If there were a smooth atlas $\mathcal B$ such that
\[
  \overline{\mathcal A}\subsetneq \mathcal B,
\]
then $\mathcal A\subset \mathcal B$, so by the previous claim we would have
\[
  \mathcal B\subset \overline{\mathcal A},
\]
a contradiction. Hence $\overline{\mathcal A}$ is maximal.

It remains to prove uniqueness.

Suppose $\mathcal B$ is another maximal smooth atlas containing $\mathcal A$.
Since $\mathcal A\subset \mathcal B$, the previous claim gives
\[
  \mathcal B\subset \overline{\mathcal A}.
\]
But $\mathcal B$ is maximal and $\overline{\mathcal A}$ is a smooth atlas. Hence
\[
  \mathcal B=\overline{\mathcal A}.
\]
Therefore the maximal smooth atlas containing $\mathcal A$ is unique.
\end{proof}

\begin{remark}
Suppose $\mathcal A_1$ and $\mathcal A_2$ are two smooth atlases on $M$ such that every chart in $\mathcal A_1$ is smoothly compatible with every chart in $\mathcal A_2$.
Then
\[
  \overline{\mathcal A_1}
  =
  \overline{\mathcal A_2}.
\]
Indeed, under this assumption,
\[
  \mathcal A_1\cup \mathcal A_2
\]
is a smooth atlas. Since it contains both $\mathcal A_1$ and $\mathcal A_2$, the uniqueness of maximal smooth atlases gives
\[
  \overline{\mathcal A_1}
  =
  \overline{\mathcal A_1\cup \mathcal A_2}
  =
  \overline{\mathcal A_2}.
\]
\end{remark}
